% Pursuing Stacks (cotes 134-2 to 134-8): archive cote 134-7, site page 99 (PDF page 100), typescript page 541.
% typescript. Draft compared: ps6.tex:4999-5057.
% AI-transcribed by Claude Opus 5.5 (claude-opus-5-5, adaptive thinking, effort high), Message Batches API, batch ps4 (msgbatch_012Aoc1NuskATPsyBeKSYK7J), 2026-11-12; not yet checked by eye.
\scanpage{134-7/99}
\mspage{541}

\margin{\unsure{a}}
a matter of chance 50/50  that we took first the choice of using W-bund-%
les,  rather than V\xout{(}-bundles which are the dual choice,  giving rise in
some respects to a simpler algebraic formalism (and to a \xout{morexintricate}
less satisfactory one in some others ...).  In both cases,  an instinct
of ``economy'' is leading us.  It is'nt always clear that instinct is'nt
misleading at times - after all,  it would be nice too to have a so-%
called ``schematic'' homotopy type (and hence a usual one) associated to
rather general types of semisimplicial schemes say.  But here already,
if we want to get a usual homotopy type by just taking sections,  we've
seen that this is'nt so automatic,  that this is tied up with the expec-%
tation that the maps (5) should be isomorphisms,  which presumably will
not be true unless we make rather drastic extra assumptions o\xout{f}n the
``bundles'' we are working with.

\margin{\unsure{.}}
Thus a first main test whether the choice of W-bundles or of
V-bundles is a work\xout{\unsure{x}}able one,  is to see whether this condition on (5)
is satisfied,  possibly under a\xout{n} \inserted[typed]{suitable} extra assumption  such as $X_{o}$=e  or
$X_{o}$=$X_{1}$=e .  The\xout{n} next test,  presumably a deeper one,  is whether these
choices allow for a description of the homotopy sheaves $\underline{\pi}_{i}(X_{*})$  in
terms of the Lie functor,under \xout{axsuitable} the natural flatness assumpti-%
on on the components of $X_{*}$  (compare  section 118).  The only clue so
far for such a relationship comes from Postnikov dévissage,  and this
relationship is'nt proved either, except in the more or less tautologi-%
cal ``Postnikov case\xout{s}'',  and without naturality.  In a way,  it appears
as a rather strange kind of relationship,  one which implies that the
homotopy type of a semisimplicial bundle is very strongly dominated
(almost determined, one might say) by \xout{thexlocal,xxxorxratherxxlocalxin}
\xout{thexsensexofxformalizationxofxschemes} it\hand{\unsure{'}}s formal completion along a
section  -  and even by the first order infinitesimal neighbourhood\xout{s}
already.  Instinct again tells us that reducing,  as it were,  a
scheme to a tangent space at one of it's points,  might make a sense
when the scheme is isomorphic to affine n-space,  and that it is surely
nonsense for general affine schemes (such as an elliptic curve minus a
point\xout{,} say !).  To say it differently,  we feel that such a thing may
be reasonable\xout{,} only when the given schemes $X_{n}$ may be thought of as
``homotopically trivial'' in some sense or other,  \xout{kinxtoxthe} which for an
algebraic cur\xout{\unsure{b}}ve say over an algebraically closed field k is surely \uline{not}
the case,  unless precisely the curve is isomorphic to the affine line.

\margin{99}
